\textbf{Purpose:} Crowd supervision can change confidence ranking without consistently improving category choice. We examine this distinction under an incomplete facial-expression output vocabulary.
\textbf{Methods:} Twenty-four matched runs on FER2013 images with FER+ votes compare five Swin-T target conditions and three ResNet-18 conditions across three seeds. Pixel-isolated selection, calibration and testing support own-ranking, fixed-image and fixed-epoch comparisons. Complete-vote disagreement is separated into unsupported mass, supported ambiguity and category-choice excess. Four adaptive phases reuse one image source.
\textbf{Results:} At 80\% coverage, selected-model Swin-T Soft-minus-Uniform disagreement is -0.38 percentage points; Soft-minus-Tie-Uniform is -0.59. ResNet-18 selected differences are +0.17 and +0.04 points; fixed-epoch differences are +0.17 and +0.19, with all four conditional intervals crossing zero. On Swin Uniform-selected images, Soft instead differs by +0.09 points. Changing the reference on identical images substantially changes F1. Soft retains more wholly unsupported images in Swin. All 36 nonempty empirical test operating points for the original four conditions exceed calibration targets.
\textbf{Conclusion:} Target representation, accepted-image composition, categorical reference and vocabulary mismatch shape apparent reliability. Cross-backbone checks extend sensitivity evidence, not independent-domain confirmation. These results support diagnostic reporting rather than universal superiority, deployment risk guarantees or validated messaging benefit.
